\chapter{الخاتمة والتوصيات}

\section{الخاتمة }


شكل هذا المشروع دراسة تحليلية شاملة للنمذجة الرياضية وتطوير استراتيجيات التحكم للمناورات الروبوتية، مطبقةً على ذراع روبوت من نوع (\LR{SCARA}). لقد انطلق مسار البحث من التوصيف الهندسي المجرد للفضاء التشغيلي عبر استنتاج معادلات الحركة المباشرة (\LR{Forward Kinematics})، وصولاً إلى بناء نموذج ديناميكي متكامل باستخدام صياغة لاغرانج-أويلر (\LR{Euler-Lagrange}). تميز هذا النموذج بدقته الفيزيائية من خلال دمج ديناميكية المحركات الكهربائية (\LR{Actuator Dynamics}) وأخذ التغيرات اللحظية في موقع مركز الثقل وعزوم القصور الذاتي للوصلة الثالثة بعين الاعتبار عند نقل الحمولات.

في شق التحكم، أثبتت الدراسة التحليلية وعمليات المحاكاة قصور خوارزميات التحكم الكلاسيكية المستقلة (\LR{PID}) عن تلبية متطلبات الأداء الصناعي عند تغير كتل التحميل، حيث أظهرت استجابة النظام انهياراً ملحوظاً في المفصل الانسحابي. لمعالجة هذه المعضلة، تم الانتقال إلى التحكم اللاخطي المتين عبر خوارزمية التحكم بالانزلاق (\LR{SMC})، والتي وفرت مناعة تامة ضد الاضطرابات، ولكن على حساب إجهاد المحركات بظاهرة الاهتزاز (\LR{Chattering}). 

توجت هذه الدراسة بتصميم واختبار خوارزمية التحكم التكيفي المعتمد على العزم المحسوب (\LR{ACTC}). أثبتت هذه الخوارزمية تفوقها الهندسي المطلق من خلال قدرتها على التكيف الذاتي، وتقدير الزيادة اللحظية في القصور الذاتي بشكل برمجي ذكي، مما أثمر عن تحقيق دقة تتبع مسار استثنائية، وزمن استقرار قياسي، وإشارات تحكم ناعمة تضمن إطالة العمر الافتراضي للمنظومة الكهروميكانيكية.

\section{التوصيات للعمل المستقبلي }

بناءً على النتائج المستخلصة من هذا البحث، وللارتقاء بأداء المنظومة الروبوتية نحو تطبيقات صناعية أكثر تعقيداً، نُدرج التوصيات الهندسية التالية للعمل المستقبلي:

\begin{enumerate}
    \item \textbf{التقدير الآلي والمسبق لكتلة الحمل ($m_L$):} 
    رغم كفاءة خوارزمية (\LR{ACTC}) في التكيف الرياضي مع القصور الذاتي المجهول، يُوصى بتزويد المؤثر النهائي للروبوت بحساسات قوة/عزم (\LR{Force/Torque Sensors})، أو تصميم مراقب حالة متقدم (\LR{State Observer}) لتقدير كتلة الحمل الديناميكية آلياً بمجرد التقاط القطعة. إدخال هذه القيمة المُقاسة مباشرةً إلى خوارزمية التحكم سيختصر زمن التكيف الابتدائي ويمنع أي انحراف أولي (\LR{Initial Undershoot}) أثناء التسارع المفاجئ.

    \item \textbf{توسيع النموذج الديناميكي للسرعات الفائقة:}
    في مهام النقل والتوضيع (\LR{Pick and Place}) التي تتطلب سرعات حركية فائقة جداً، يُوصى بإعادة دمج مصفوفة قوى كوريوليس والطرد المركزي $C(q, \dot{q})$ التي تم إهمالها في هذه الدراسة. سيؤدي دمج هذه القوى غير الخطية إلى زيادة التعقيد الحسابي، ولكنه سيرفع من دقة النموذج التنبؤي في ظروف التشغيل القاسية.
\end{enumerate}